\documentclass[12pt,a4paper]{article}

\usepackage{cmap}
\usepackage[T2A]{fontenc}
\usepackage[utf8]{inputenc}
\usepackage[russian]{babel}
\usepackage{amsmath,amssymb}
\usepackage[a4paper,margin=2cm]{geometry}

\setlength{\parindent}{0pt}
\setlength{\parskip}{0.45em}
\sloppy

\begin{document}

\begin{center}
{\Large\bfseries Билет 2: вопросы для проверки знаний}
\end{center}

\noindent\fbox{%
  \begin{minipage}{\dimexpr\linewidth-2\fboxsep-2\fboxrule\relax}
  \normalfont
Предел числовой функции нескольких переменных. Предел функции по множеству.
Пределы по направлениям. Непрерывность функции нескольких переменных в точке и
по множеству. Непрерывность сложной функции. Свойства функций, непрерывных на
компакте: ограниченность, достижимость точных нижней и верхней граней,
равномерная непрерывность. Теорема о промежуточных значениях функции,
непрерывной в области.
  \end{minipage}%
}

\section*{Вопросы на понимание}

\begin{enumerate}
\item Чем предел функции по совокупности переменных отличается от предела по множеству?
\item Почему в определении предела по множеству важно, чтобы точка была предельной точкой этого множества?
\item Что меняется в непрерывности функции по множеству, если рассматриваемая точка является изолированной?
\item Если общий предел функции в точке существует, что можно сказать о пределах по направлениям? Почему обратное рассуждение опасно?
\item Почему совпадение пределов по всем прямым направлениям еще не гарантирует существование общего предела?
\item Как с помощью двух разных кривых, подходящих к одной точке, можно опровергнуть существование предела функции двух переменных?
\item Чем обычная непрерывность в каждой точке множества отличается от равномерной непрерывности на этом множестве?
\item Почему в теоремах об ограниченности и достижении точных граней нельзя просто заменить компакт произвольным множеством?
\item Приведите пример непрерывной функции на некомпактном множестве, которая ограничена, но не достигает точной верхней или нижней грани.
\item Как непрерывность внутренней и внешней функций вместе обеспечивает непрерывность композиции?
\item Почему в теореме о промежуточных значениях для функции нескольких переменных нужно условие связности области?
\item Как идея соединения двух точек кривой сводит теорему о промежуточных значениях в области к одномерному случаю?
\end{enumerate}

\section*{Источники для сверки}

\texttt{target/answer\_question\_02.tex};
\texttt{IvGE\_1.pdf}, стр. 57--60, 167--179;
\texttt{IvGE\_2.pdf}, стр. 57.

\end{document}
